%% Lec03 shared MAP slide: "one map, four acts" recurring diagram.
%% Input this file AFTER ../shared_preamble.tex (needs RedTitle, VeryLightGray,
%% DarkText) and AFTER \usetikzlibrary{positioning,arrows.meta,shapes,calc}.
%% Call \LecThreeMap{k} inside a frame, where k in {1,2,3,4} highlights the
%% current act ("you are here") and k = 0 renders the closing reprise with all
%% four acts checked green (used by T4's closing frame).
%% Filename deliberately does NOT match the build glob Lec*_T*.tex, so the build
%% script never tries to compile it standalone.
%%
%% The four acts are the lecture's story: WHY -> BUILD -> TRAIN -> SHAPE, i.e.
%% the case for ML in economics, the network as an object (architecture + UAT),
%% the learning machinery (backprop + SGD + regularization + double descent),
%% and economics-aware architectures with the hands-on lab.

\newcommand{\LecThreeMap}[1]{%
% Precompute per-box style selectors; TikZ option lists are comma-split before
% TeX conditionals run, so \ifnum must not sit inside a [...] options list.
\ifnum#1=0
  \tikzset{lmapa/.style={lmapdone}, lmapb/.style={lmapdone},
           lmapc/.style={lmapdone}, lmapd/.style={lmapdone}}%
\else
  \tikzset{lmapa/.style={}, lmapb/.style={}, lmapc/.style={}, lmapd/.style={}}%
  \ifnum#1=1 \tikzset{lmapa/.style={lmapcur}}\fi
  \ifnum#1=2 \tikzset{lmapb/.style={lmapcur}}\fi
  \ifnum#1=3 \tikzset{lmapc/.style={lmapcur}}\fi
  \ifnum#1=4 \tikzset{lmapd/.style={lmapcur}}\fi
\fi
\begin{center}
\begin{tikzpicture}[
    act/.style={rectangle, rounded corners, draw=black!60, fill=VeryLightGray,
        minimum width=3.0cm, minimum height=0.95cm, text width=2.9cm,
        align=center, font=\small, text=DarkText},
    lmapcur/.style={draw=RedTitle, very thick, fill=orange!20},
    lmapdone/.style={draw=green!45!black, thick, fill=green!10},
    q/.style={font=\tiny\itshape, text width=3.0cm, align=center, text=black!70},
    barlab/.style={font=\tiny\bfseries, text=black!65},
    sat/.style={rectangle, rounded corners, draw=black!45, dashed, fill=white,
        text width=5.2cm, align=center, font=\tiny, text=black!70,
        minimum height=0.5cm},
    arr/.style={-{Stealth[length=2.2mm]}, thick, gray!60}
]
    % --- four act boxes ---
    \node[act, lmapa] (a1) at (0,0)
        {\textbf{3.1\ \ WHY}\\[1pt] \tiny the case for ML in econ};
    \node[act, lmapb] (a2) at (3.6,0)
        {\textbf{3.2\ \ BUILD}\\[1pt] \tiny neurons $\cdot$ activations $\cdot$ UAT};
    \node[act, lmapc] (a3) at (7.2,0)
        {\textbf{3.3\ \ TRAIN}\\[1pt] \tiny backprop $\cdot$ SGD $\cdot$ dbl.\ descent};
    \node[act, lmapd] (a4) at (10.8,0)
        {\textbf{3.4\ \ SHAPE}\\[1pt] \tiny ICNN $\cdot$ monotone $\cdot$ lab};
    \draw[arr] (a1) -- (a2);
    \draw[arr] (a2) -- (a3);
    \draw[arr] (a3) -- (a4);
    % --- the student question each act answers ---
    \node[q] at (0,-0.95)    {why should an economist\\ care about ML?};
    \node[q] at (3.6,-0.95)  {what is a neural network,\\ and why does it work?};
    \node[q] at (7.2,-0.95)  {how does it learn ---\\ and why not overfit?};
    \node[q] at (10.8,-0.95) {how do I build economics\\ in --- and run it myself?};
    % --- "you are here" marker above the current act ---
    \ifnum#1=1 \node[font=\scriptsize\bfseries, text=RedTitle] at (0,0.98)
        {you are here\ $\blacktriangledown$};\fi
    \ifnum#1=2 \node[font=\scriptsize\bfseries, text=RedTitle] at (3.6,0.98)
        {you are here\ $\blacktriangledown$};\fi
    \ifnum#1=3 \node[font=\scriptsize\bfseries, text=RedTitle] at (7.2,0.98)
        {you are here\ $\blacktriangledown$};\fi
    \ifnum#1=4 \node[font=\scriptsize\bfseries, text=RedTitle] at (10.8,0.98)
        {you are here\ $\blacktriangledown$};\fi
    \ifnum#1=0 \node[font=\scriptsize\bfseries, text=green!45!black] at (5.4,0.98)
        {all four acts complete};\fi
    % --- bottom band: the ML pipeline (the technical spine of the lecture) ---
    \fill[blue!10]   (-1.55,-1.72) rectangle (5.15,-1.40);
    \fill[green!12]  (5.65,-1.72)  rectangle (8.75,-1.40);
    \fill[orange!16] (9.25,-1.72)  rectangle (12.35,-1.40);
    \node[barlab] at (1.8,-1.56)  {REPRESENT\ (UAT $\cdot$ Barron)};
    \node[barlab] at (7.2,-1.56)  {OPTIMIZE\ (loss $\cdot$ gradient)};
    \node[barlab] at (10.8,-1.56) {GENERALIZE $+$ CONSTRAIN};
    % --- dashed satellites: where this lecture sits in the course flow ---
    \node[sat] at (1.5,-2.42)
        {\textit{upstream:} Lec01--02 --- AI as cognitive capital; deep learning as the engine of the predictive stack};
    \node[sat] at (9.3,-2.42)
        {\textit{downstream:} Lec04--06 --- these networks become Bellman/Euler solvers for dynamic macro models};
\end{tikzpicture}
\end{center}
}
